% Source problem: https://www.ap.t.u-tokyo.ac.jp/themes/ap_t_u-tokyo/graduate/pdf/2027ap_exam_je.pdf
\documentclass[11pt,a4paper]{article}
\usepackage[margin=24mm]{geometry}
\usepackage{fontspec}
\usepackage{xeCJK}
\setCJKmainfont{HaranoAjiMincho-Regular.otf}
\setCJKsansfont{HaranoAjiGothic-Medium.otf}
\usepackage{amsmath,amssymb}
\usepackage{tikz}
\setlength{\parindent}{1em}
\setlength{\parskip}{0.35em}
\title{物理工学専攻 院試解答集\\2027年度 物理学 第4問}
\author{政岡凜太郎}
\date{}
\begin{document}
\maketitle

\section*{[1.1] 振動子の応答}
電場方向の単位ベクトルを $\boldsymbol e_x$ とする。$\boldsymbol u(t)=u_0\boldsymbol e_xe^{-i\omega t}$ を (1) に代入すると、$\mathrm d^2\boldsymbol u/\mathrm dt^2=-\omega^2\boldsymbol u$ より
\begin{equation}
m(\omega_0^2-\omega^2)u_0=qE_0,
\qquad
\boxed{P_0=Nqu_0=\frac{Nq^2E_0}{m(\omega_0^2-\omega^2)}}.
\end{equation}
この定常調和応答は $\omega\ne\omega_0$ に限る。減衰のないモデルでちょうど $\omega=\omega_0$ とすると振幅が有限な同じ周波数の定常解はなく、式の発散をそのまま有限の分極と解釈できない。

\section*{[1.2] 誘電率とその符号}
以後、振動子の強さを周波数の二乗で表した
\begin{equation}
\Omega_p^2\equiv\frac{Nq^2}{\varepsilon_0m}>0
\end{equation}
を使う。[1.1] のベクトル関係 $\boldsymbol P=\varepsilon_0\Omega_p^2\boldsymbol E/(\omega_0^2-\omega^2)$ を誘電率の定義に入れると
\begin{equation}
\boxed{\varepsilon(\omega)=\varepsilon_0\left(1+\frac{\Omega_p^2}{\omega_0^2-\omega^2}\right)
=\varepsilon_0\frac{\omega_0^2+\Omega_p^2-\omega^2}{\omega_0^2-\omega^2}}.
\end{equation}
$\omega>0$ で分母と分子の符号が異なるのは
\begin{equation}
\boxed{\omega_0<\omega<\sqrt{\omega_0^2+\Omega_p^2}}
\end{equation}
である。下端は共鳴点、上端は $\varepsilon=0$ なので、いずれも「負」には含まない。

\section*{[1.3] 誘電体中の固有モード}
(4) の回転をとり (3) を代入すると
\begin{equation}
\nabla\times(\nabla\times\boldsymbol E)
=i\omega\mu_0\nabla\times\boldsymbol H
=\mu_0\omega^2(\varepsilon_0\boldsymbol E+\boldsymbol P).
\label{eq:double-curl}
\end{equation}
一様な誘電体で $\varepsilon(\omega)\ne0$ なら、(5) の
$\nabla\cdot\boldsymbol D=\varepsilon(\omega)\nabla\cdot\boldsymbol E=0$ から $\boldsymbol k\cdot\boldsymbol E_0=0$。したがって (2) の平面波について二重回転の左辺は $k^2\boldsymbol E$ となる。非零の横電場では
\begin{equation}
\boxed{\boldsymbol P=\left(\frac{k^2}{\mu_0\omega^2}-\varepsilon_0\right)\boldsymbol E},
\qquad
\boxed{c^2k^2=\omega^2\frac{\varepsilon(\omega)}{\varepsilon_0}
=\omega^2\left(1+\frac{\Omega_p^2}{\omega_0^2-\omega^2}\right)}.
\label{eq:bulk}
\end{equation}
最初の式はマクスウェル方程式による分極と電場の関係であり、[1.1] の振動子応答を加えたものが二番目の分散式である。

$K=c^2k^2$、$X=\omega^2$ と置き、分母を払うと
\begin{equation}
X^2-(K+\omega_0^2+\Omega_p^2)X+K\omega_0^2=0.
\end{equation}
ゆえに、横波の二枝は
\begin{equation}
\boxed{\omega_{\pm}^2(k)=\frac{K+\omega_0^2+\Omega_p^2
\pm\sqrt{(K+\omega_0^2+\Omega_p^2)^2-4K\omega_0^2}}{2}}.
\label{eq:branches}
\end{equation}
判別式は $(K-\omega_0^2)^2+2\Omega_p^2(K+\omega_0^2)+\Omega_p^4>0$ で、$k>0$ なら二根は正である。$k=0$ では下枝は $\omega_-=0$、上枝は $\sqrt{\omega_0^2+\Omega_p^2}$ となる。

なお $\varepsilon(\omega)=0$ では (5) から横波条件を導けない。このとき $\omega=\sqrt{\omega_0^2+\Omega_p^2}$、$\boldsymbol E\parallel\boldsymbol k$、$\boldsymbol H=\boldsymbol 0$ の縦の分極振動も (3)--(6) を満たす。式 \eqref{eq:branches} は非零横電場の分散を表す。

\section*{[1.4] 大きな波数での二枝}
$K\gg\omega_0^2,\Omega_p^2$ では、\eqref{eq:branches} の根の和と積から先頭項が分かる。設問で指定された三つの記号で表す答えは
\begin{equation}
\boxed{\omega_+\simeq ck,\qquad\omega_-\simeq\omega_0}.
\end{equation}
根の和と積を順に展開すると
\begin{equation}
\omega_+^2=K+\Omega_p^2+O(K^{-1}),\qquad
\omega_-^2=\omega_0^2\left(1-\frac{\Omega_p^2}{K}+O(K^{-2})\right).
\end{equation}
平方根を一次まで展開した結果は
\begin{equation}
\omega_+=ck+\frac{\Omega_p^2}{2ck}+O(k^{-3}),\qquad
\omega_-=\omega_0\left(1-\frac{\Omega_p^2}{2c^2k^2}+O(k^{-4})\right).
\label{eq:highk}
\end{equation}
この次項の係数には $\Omega_p^2=Nq^2/(\varepsilon_0m)$ が現れる。$ck$ と $\omega_0$ はいずれも角周波数の次元を持ち、$k\to\infty$ で補正はそれぞれの先頭項に比べて消える。

\section*{[2.1] 界面の連続条件と電場成分}
図のように界面の上を真空、下を誘電体とし、$\alpha_0,\alpha_1$ は局在のため正の実数とする。
\begin{center}
\begin{tikzpicture}[>=stealth,scale=0.85]
  \fill[gray!12] (-3,-1.25) rectangle (3,0);
  \draw[thick] (-3,0)--(3,0);
  \draw[->] (0,-1.5)--(0,1.7) node[above] {$z$};
  \draw[->] (-3.2,0)--(3.35,0) node[right] {$x$};
  \node at (-2.1,0.9) {真空 $\varepsilon_0$};
  \node at (-2.1,-0.8) {誘電体 $\varepsilon(\omega)$};
  \draw[->,very thick] (0.6,0.23)--(2.0,0.23) node[midway,above] {$k_{\parallel}$};
  \draw[->,blue,thick] (2.45,0.15)--(2.45,1.15) node[above] {$e^{-\alpha_0 z}$};
  \draw[->,blue,thick] (2.45,-0.15)--(2.45,-1.15) node[below] {$e^{\alpha_1 z}$};
\end{tikzpicture}
\end{center}
\newpage
\enlargethispage{3\baselineskip}
[1] の無源 Maxwell 方程式を界面をまたいで用いる。(4) を界面を横切る微小ループで積分すると接線電場が連続し、(5) を薄い円柱で積分すると自由表面電荷がない場合の法線電束密度が連続する。場が界面近傍で有限なら、縮小極限でループ短辺の線積分、ループの磁束、円柱側面の電束は消える。$y$ 成分は (7), (8) により零だから、$z=0$ では
\begin{equation}
\boxed{E_{0x}=E_{1x}},\qquad
\boxed{\varepsilon_0E_{0z}=\varepsilon(\omega)E_{1z}}.
\label{eq:boundary}
\end{equation}
真空側では $\varepsilon_0\ne0$ なので、(5) と (7) から $ik_{\parallel}E_{0x}-\alpha_0E_{0z}=0$ を得る。媒質側では $\varepsilon=0$ の可能性を先に除かず、(3), (4) から得る二重回転方程式の $x$ 成分を用いる。(8) と $\alpha_1^2=k_{\parallel}^2-\varepsilon\mu_0\omega^2$ を代入すると
\begin{equation}
-\alpha_1(\alpha_1E_{1x}-ik_{\parallel}E_{1z})
=\varepsilon\mu_0\omega^2E_{1x}
\quad\Longrightarrow\quad
ik_{\parallel}E_{1x}+\alpha_1E_{1z}=0.
\end{equation}
局在条件 $\alpha_0>0$ から $k_{\parallel}^2>\varepsilon_0\mu_0\omega^2$ であり、$k_{\parallel}\ne0$ である。従って
\begin{equation}
\boxed{E_{0z}=\frac{ik_{\parallel}}{\alpha_0}E_{0x}},\qquad
\boxed{E_{1z}=-\frac{ik_{\parallel}}{\alpha_1}E_{1x}}.
\label{eq:component}
\end{equation}
法線成分の符号が逆なのは、両側で指数の $z$ 微分の符号が逆だからである。
$\varepsilon=0$ なら式\eqref{eq:boundary}の法線電束の連続性から $E_{0z}=0$、式\eqref{eq:component}と接線電場の連続性から他の電場成分も零となる。従って (7), (8) の形の非零の局在表面波ではこの場合を除ける。

\section*{[2.2] 表面波の条件}
\eqref{eq:component} を \eqref{eq:boundary} に入れ、非零の $E_{0x}=E_{1x}$ と $k_{\parallel}$ を消すと
\begin{equation}
\boxed{\frac{\varepsilon_0}{\alpha_0}+\frac{\varepsilon(\omega)}{\alpha_1}=0,
\quad\text{すなわち}\quad
\varepsilon(\omega)=-\varepsilon_0\frac{\alpha_1}{\alpha_0}}.
\label{eq:tm}
\end{equation}
$\alpha_0,\alpha_1>0$ なので少なくとも $\varepsilon(\omega)<0$ が必要である。実際の存在範囲は両側の減衰条件を合わせて [2.4] で決める。

\section*{[2.3] 界面に沿う波数}
\eqref{eq:tm} より $\alpha_1=-\varepsilon\alpha_0/\varepsilon_0$。両辺の二乗と
$\alpha_0^2=k_{\parallel}^2-\varepsilon_0\mu_0\omega^2$、
$\alpha_1^2=k_{\parallel}^2-\varepsilon\mu_0\omega^2$ を組み合わせると
\begin{equation}
(\varepsilon_0^2-\varepsilon^2)k_{\parallel}^2
=\varepsilon_0\varepsilon\mu_0\omega^2(\varepsilon_0-\varepsilon).
\end{equation}
$\varepsilon\ne\varepsilon_0$ として共通因子を除き、正方向へ伝わる根を選べば
\begin{equation}
\boxed{k_{\parallel}=\omega\sqrt{\frac{\mu_0\varepsilon_0\varepsilon(\omega)}
{\varepsilon_0+\varepsilon(\omega)}}
=\frac{\omega}{c}\sqrt{\frac{\varepsilon(\omega)}
{\varepsilon_0+\varepsilon(\omega)}}}.
\label{eq:kpar}
\end{equation}
これは $\varepsilon<-\varepsilon_0$ の範囲で実数かつ $k_{\parallel}>\omega/c$ となり、真空側でも減衰する。

\section*{[2.4] 表面波が存在する周波数}
\eqref{eq:kpar} を $\alpha_0^2=k_{\parallel}^2-\omega^2/c^2$ に入れると
\begin{equation}
\alpha_0^2=\frac{\omega^2}{c^2}\frac{-\varepsilon_0}{\varepsilon+\varepsilon_0}.
\end{equation}
$\alpha_0^2>0$ と [2.2] の $\varepsilon<0$ を両立する条件は $\varepsilon<-\varepsilon_0$。このとき $\alpha_1=-\varepsilon\alpha_0/\varepsilon_0>0$ も満たす。[1.2] の誘電率を代入すると
\begin{equation}
1+\frac{\Omega_p^2}{\omega_0^2-\omega^2}<-1
\quad\Longleftrightarrow\quad
\boxed{\omega_0<\omega<\sqrt{\omega_0^2+\frac{Nq^2}{2\varepsilon_0m}}}.
\end{equation}
下端 $\omega_0$ では無減衰振動子の定常応答が発散するうえ、真空側の $\alpha_0\to0$ で局在性も失う。上端では $\varepsilon\to-\varepsilon_0$ に伴い $k_{\parallel}$ が無限大になる。よって両端を除く。[1.2] の $\varepsilon<0$ の帯域のうち、この狭い範囲だけが真空・誘電体の両側へ減衰する表面波を与える。
\end{document}
